%!TeX root=../thesis.tex

\chapter{Discussion}
\label{cap:discussion}

% TODO: interpret the results and relate them to the research questions.
% Suggested structure (adapt as needed):
%
% \section{Answering the Research Questions}
% Revisit each RQ from Chapter 1 and discuss how the findings address it.
%
% \section{Comparison with Prior ArchHypo Evaluations}
% Compare your findings with the TSE study (Jetsoft/PropSEP) and the
% IEEE Software study (Catch Solve). What is similar? What is different?
% What does tool support change?
%
% \section{Implications for Practice}
% What do the results mean for teams considering ArchHypo/Hypo-Stage?
% What recommendations can you offer?
%
% \section{Implications for Research}
% What new questions emerged? What should future work investigate?
%
% \section{Threats to Validity}
% Discuss internal, external, construct, and conclusion validity threats.
% Be specific: what biases exist? How did you mitigate them?
%
% TIP: use \textbf{Finding N} cross-references to tie discussion back to results.
